\setlength{\footskip}{8mm}

\chapter{EXPERIMENTAL PLAN AND PROJECT MANAGEMENT}

\section{Baseline and Model Variants}

The experimental design compares SSM-Adapter against controlled baselines on two independent backbones. The most important fairness requirement is that the parameter-matched transformer control variant must use the same total added parameter count as SSM-Adapter and be inserted at the same point in the pipeline. This prevents the experiment from confusing the benefit of state-space temporal integration with the benefit of simply adding more parameters at that location.

\begin{table}[h]
\caption{Planned model variants across two backbones.}
\begin{center}
\begin{tabular}{p{0.55in}p{2.8in}p{2.3in}}
\toprule
ID & Model variant & Purpose \\
\midrule
E0 & pgat-length v2 baseline (mBART, no adapter) & Same-backbone lower reference. \\
E1 & pgat-length v2 with SSM-Adapter & Main pgat-length experimental system. \\
E2 & pgat-length v2 with parameter-matched transformer adapter & Controls for added parameters vs SSM specifically. \\
E3 & pgat-length v2 with forward-only SSM (no backward direction) & Ablates bidirectionality. \\
E4 & TSPNet baseline reproduction (I3D + transformer, no adapter) & Second-backbone lower reference. \\
E5 & TSPNet with SSM-Adapter & Main TSPNet experimental system. \\
E6 & TSPNet with parameter-matched transformer adapter & Controls for added parameters on the second backbone. \\
\bottomrule
\end{tabular}
\label{tab:model-variants}
\end{center}
\end{table}

The primary claim is supported when E1 improves over E0 and E5 improves over E4 on the length-cliff dimension by an amount not explained by E2 relative to E0 or E6 relative to E4. The claim is falsified if either backbone shows no improvement or if the parameter-matched transformer control produces a comparable improvement.

\section{Ablation Studies}

The ablation plan is designed to answer which parts of the SSM design actually matter. Bidirectionality is tested by comparing E3 with E1. State dimension is tested by sweeping $d_{\text{state}} \in \{8, 16, 32, 64\}$ on pgat-length. Layer count per direction is tested by sweeping $L \in \{1, 2, 3\}$. Residual scaling is tested at $s \in \{0.5, 1.0, 2.0\}$. The compute budget for each ablation is one training run per hyperparameter setting on pgat-length.

\begin{table}[h]
\caption{Ablation comparisons and interpretation.}
\begin{center}
\begin{tabular}{p{1.4in}p{2.1in}p{2.0in}}
\toprule
Ablation & Comparison & Interpretation \\
\midrule
Adapter presence & E1 vs E0, E5 vs E4 & Primary effect of the module. \\
Parameter matching & E1 vs E2, E5 vs E6 & Whether SSM structure matters beyond added parameters. \\
Bidirectionality & E1 vs E3 & Whether both directions are needed. \\
State dim & $d_{\text{state}} \in \{8,16,32,64\}$ & Capacity-quality trade-off. \\
Layer count & $L \in \{1,2,3\}$ per direction & Depth-quality trade-off. \\
Residual scale & $s \in \{0.5,1.0,2.0\}$ & Sensitivity to residual weighting. \\
Cross-backbone & Consistency of length-cliff reduction across pgat-length and TSPNet & Backbone-agnostic claim. \\
\bottomrule
\end{tabular}
\label{tab:ablations}
\end{center}
\end{table}

\section{Metrics and Statistical Reporting}

chrF is a primary metric because it captures character n-gram overlap and is suitable for German morphological variation \citep{popovic2015chrf}. BLEU-1, BLEU-2, BLEU-3, and BLEU-4 are reported for comparability with prior SLT work, using SacreBLEU signatures for reproducibility \citep{post2018sacrebleu}. ROUGE-L is reported as a supporting sequence-overlap metric \citep{lin2004rouge}. All external baselines are rescored under the same scorer so cross-system comparisons are apples-to-apples.

The primary evaluation is a five-bin length-stratified analysis. Test samples are grouped by whitespace-tokenised reference word count into five bins (1--6, 7--12, 13--18, 19--24, 25--32 words) and each metric is reported per bin along with the overall corpus number. A single-number summary of the length cliff is the long-over-short ratio, defined as the mean metric on bins 13--32 divided by the mean metric on bins 1--12. A higher ratio indicates a flatter length curve.

All main results are reported as mean and standard deviation across three random seeds where compute allows. Statistical significance is assessed using paired bootstrap resampling with 2000 samples for the main comparisons: E1 versus E0, E1 versus E2, E5 versus E4, and E5 versus E6. Reported deltas include 95 percent bootstrap confidence intervals.

\begin{table}[h]
\caption{Evaluation metrics and reporting plan.}
\begin{center}
\begin{tabular}{p{1.2in}p{2.05in}p{2.1in}}
\toprule
Metric or analysis & Purpose & Reporting plan \\
\midrule
BLEU-4 & Standard machine translation metric for prior-work comparability. & Overall and per-bin, SacreBLEU 13a, exponential smoothing. \\
chrF & Character-level metric well-suited to German morphology. & Overall and per-bin, char\_order 6, $\beta = 2$. \\
ROUGE-L F1 & Sequence-overlap metric based on longest common subsequence. & Overall and per-bin, no stemmer. \\
Five-bin length analysis & Isolates the length cliff and measures whether the adapter flattens it. & Bins 1--6, 7--12, 13--18, 19--24, 25--32; long-over-short ratio reported. \\
Bootstrap significance & Tests whether adapter deltas are reliable. & Paired bootstrap with 2000 resamples on E1 vs E0, E1 vs E2, E5 vs E4, E5 vs E6. \\
Cross-backbone consistency & Tests whether the effect transfers. & Direction and magnitude of the length-cliff delta compared across pgat-length and TSPNet. \\
Prediction inspection & Qualitative check for mode collapse. & Manual review of 20 randomly sampled predictions per system focusing on distinct outputs across distinct videos. \\
\bottomrule
\end{tabular}
\label{tab:metrics-reporting}
\end{center}
\end{table}

\section{Feasibility and Compute Plan}

The method is designed for a single 48 GB GPU (A6000). pgat-length training uses cached DINOv2 and TimeSformer features so no visual backbone runs during training; mBART fine-tune with a variable-length prefix and the SSM adapter fits in bf16 with gradient checkpointing at micro-batch three, gradient accumulation eight. TSPNet training uses pre-extracted I3D features and a small transformer decoder, easily fitting on the same GPU at max-tokens 4096.

Wall-time estimates per training run are as follows. On pgat-length, one full run is approximately five to eight hours until convergence via early stopping. On TSPNet, one full run is approximately four to eight hours to converge, with early stopping via reduce-on-plateau. The full experimental plan (E0 to E6 plus ablations across three seeds where compute allows) fits within roughly forty A6000-days.

Storage is bounded by cached features and checkpoints. PHOENIX-2014T raw frames and cached DINOv2/TimeSformer features use approximately fifty-five gigabytes; TSPNet's downloaded I3D features and BPE data-bin use approximately eight gigabytes after extraction; the required Hugging Face model cache uses approximately six gigabytes with the Qwen entry removed and mBART retained. Adapter checkpoints are small (five to twenty megabytes) and add negligible storage. The full storage budget fits within a one-hundred gigabyte quota with margin.

\section{Risk Management}

The first major risk is that the adapter improves one backbone but not the other, weakening the cross-backbone claim. The mitigation is to report the negative result honestly and treat the finding as backbone-specific rather than backbone-agnostic; the release then documents the adapter as a candidate module and the paper narrative reduces its scope.

The second risk is TSPNet reproducibility. The published TSPNet code is a fairseq fork requiring old PyTorch and old NumPy pins. The mitigation is to install TSPNet in a separate conda environment with pinned versions (PyTorch 1.11, NumPy 1.19), apply small surgical patches for known PyTorch 1.11 API changes (integer floor division in the beam-search buffer), and use sacrebleu 1.5.1 to match the fairseq validation code. The successful local reproduction confirms this path is viable.

The third risk is that both backbones' improvements are small enough that the primary claim relies on chrF and ROUGE-L flattening rather than BLEU-4 improvement. The mitigation is to report all three metrics from the start, discuss the well-documented lag of BLEU-4 behind chrF on morphologically rich languages such as German, and commit to a specific plan for improving BLEU-4 through downstream techniques (length-penalty tuning, sequence-level knowledge distillation, or coverage penalties at inference).

\section{Decision Rules and Contingency Plan}

\begin{itemize}
    \item Milestone 1: pgat-length baseline (E0) and pgat-length with SSM adapter (E1) both complete on the development split. If the length-cliff delta between E1 and E0 is not directional (flatter curve), the study pauses to inspect predictions for mode collapse before proceeding to the second backbone.
    \item Milestone 2: TSPNet baseline (E4) reproduces the published BLEU-4 near $12.98$ under the local scorer within one BLEU-4 point. If it does not, the reproduction protocol is examined before running the adapter variant on TSPNet.
    \item Milestone 3: TSPNet with SSM adapter (E5) completes. If E5 improves length-cliff metrics over E4, the primary claim is supported; if it does not, the study is honestly reported as a backbone-specific improvement.
    \item Milestone 4: Parameter-matched transformer adapter controls (E2 and E6) complete. If E2 matches E1 and E6 matches E5, the SSM-specific claim is weakened and the discussion attributes the effect to the added-parameter position rather than the state-space dynamics.
    \item Milestone 5: Ablations complete on state dim, layers, direction, and residual scale. If no ablation shifts the length-cliff metric meaningfully, the adapter's effect is reported as robust to hyperparameter choice.
    \item Milestone 6: The final frozen-architecture PHOENIX-2014T test run is executed once on each system for the paper. Test scores are not used for design decisions.
\end{itemize}

\section{Thesis Timeline}

The revised thesis timeline runs from September 2026 to May 2027, with the SSM-Adapter experimental programme scheduled as the primary experimental contribution.

\begin{landscape}
\begin{figure}[h]
\caption{Thesis timeline from September 2026 to May 2027.}
\centerline{\includegraphics[width=8.9in]{figures/timeline_placeholder.png}}
\label{fig:research-timeline}
\small{\textit{Note.} The timeline covers baseline replication and pgat-length adapter runs first, then TSPNet reproduction and adapter runs, then ablations and cross-backbone analysis, then writing and final revision.}
\end{figure}
\end{landscape}
